% Generated draft: requires mathematical and source review.
\documentclass[11pt]{article}
\usepackage[margin=25mm]{geometry}
\usepackage{fontspec}
\setmainfont{Latin Modern Roman}
\usepackage{amsmath,amssymb,mathtools,unicode-math}
\setmathfont{Latin Modern Math}
\usepackage{enumitem,needspace,xurl,hyperref}
\setlength{\parindent}{0pt}
\setlength{\parskip}{5pt}
\setlength{\emergencystretch}{3em}
\newcommand{\N}{\mathbb{N}}
\newcommand{\Z}{\mathbb{Z}}
\newcommand{\Q}{\mathbb{Q}}
\newcommand{\R}{\mathbb{R}}
\newcommand{\C}{\mathbb{C}}
\newcommand{\F}{\mathbb{F}}
\newcommand{\A}{\mathbb{A}}
\newcommand{\PP}{\mathbb{P}}
\newcommand{\E}{\mathbb{E}}
\newcommand{\eps}{\varepsilon}
\newcommand{\dd}{\,\mathrm{d}}
\begin{document}
\section*{UnsolvedMath ID: 30004247}
\textbf{Source problem code:} catalogue entry
\subsection{Definitions and mathematical statement}
The supplied question concerns a kinetically constrained model (KCM), an initial infection law \(\mu\), and an associated bootstrap-percolation process. The underlying sites, configuration space, kinetic constraints, transition rates, bootstrap update rule, and boundary conditions are not specified in the record; these must be supplied before a precise comparison can be formulated.

For any such specified model with a distinguished site denoted by \(0\), let \(\eta_t\) denote the KCM configuration at time \(t\), started with distribution \(\mu\). Define its infection hitting time by
\[
\tau_0=\inf\{t\geq 0: \text{site }0\text{ is infected in }\eta_t\},
\qquad \inf\varnothing=+\infty.
\]
Then \(\mathbb E_\mu[\tau_0]\) averages over the initial configuration and the subsequent KCM evolution, and may be infinite. The record does not specify whether any additional random environment is present or included in this expectation.

Write \(T_{\mathrm{rel}}\) for the KCM relaxation time. Under the conventional interpretation for a reversible continuous-time Markov process with generator \(L\) and invariant probability measure \(\pi\), this means
\[
T_{\mathrm{rel}}=\bigl(\operatorname{gap}_{L^2(\pi)}(-L)\bigr)^{-1},
\]
with value \(+\infty\) when the gap is zero. This is a conditional clarification of terminology: the source record does not supply \(L\), \(\pi\), reversibility, or a rate normalization, and does not identify the initial law \(\mu\) with an invariant law.

Let \(\tau_{\mathrm{BP}}\) denote the bootstrap-percolation infection time mentioned in the source, with initial configuration distributed according to \(\mu\). Its target site or target set, time convention, and precise relation to the KCM constraints remain unspecified. In particular, the record does not define which statistic of its distribution is meant by its ``typical value.''

\textbf{Question as supplied.} Determine how the asymptotic scaling of \(T_{\mathrm{rel}}\) and \(\mathbb E_\mu[\tau_0]\) is related to the typical bootstrap-percolation infection time under \(\mu\). The edited statement asks more specifically whether these quantities have the same asymptotic scale, whereas the original statement asks whether their scaling is connected. These formulations have different strengths and are not identified here.

A precise asymptotic assertion requires a specified family of models and initial laws, a limiting parameter and regime, a definition of ``typical,'' and a comparison relation. For example, bounds up to multiplicative constants and equality of leading logarithmic asymptotics are distinct possible meanings of ``same asymptotic scale''; the record does not select either. It also does not state the admissible class of laws \(\mu\), the hypotheses ensuring finite infection times, or whether comparison constants must be uniform in \(\mu\). Consequently, the supplied evidence supports this qualified question, but does not determine a unique quantified conjecture.
\subsection{Short English statement}
For a kinetically constrained model started from an infection law \(\mu\), ask how its relaxation time and expected infection time are related to a typical infection time in associated bootstrap percolation. The model, limiting regime, admissible initial laws, meaning of ``typical,'' and strength of the proposed comparison are unspecified. The original question asks for a connection; the edited version asks for the same asymptotic scale.
\subsection{Sources}
\begin{itemize}
\item Local registry: \texttt{data/unsolvedmath/problems.json}, permanent UnsolvedMath ID \(30004247\), source problem code \texttt{catalogue entry}, title ``Kinetic Relaxation and Bootstrap Percolation Infection Times.'' The supplied record belongs to the Oberwolfach Reports Open Problems collection and gives proposed year 2019; no individual proposer is identified.
\item Supplied source citation: \emph{Large Scale Stochastic Dynamics} (2020), \emph{Oberwolfach Reports}, pp.~2605--2669, DOI 10.4171/owr/2019/42. \url{https://doi.org/10.4171/owr/2019/42}. The record describes the report as documenting workshop discussions and open problems. No problem-specific page or speaker attribution is supplied.
\item Supplied literature source, labeled ``Random-constraint KCM paper'': \url{https://arxiv.org/abs/1812.00774}. No authors or full bibliographic citation are supplied in the record.
\item Supplied literature source, labeled ``Refined universality paper'': \url{https://doi.org/10.1007/s00220-023-04874-8}. No authors or full bibliographic citation are supplied in the record.
\item Literature assessment checked at \texttt{2026-08-22}: the supplied dated curation labels the problem ``partially solved'' and states that random-constraint KCM work establishes direct comparisons and refined critical scaling bounds involving bootstrap-percolation infection times, while the full proposed equivalence for arbitrary laws \(\mu\) is not proved. This is the supplied assessment derived from OWR curation, not an independently verified literature conclusion; its status and cited claims require independent checking.
\item The registry labels the statement \texttt{corrected\_verified} and describes repairs to notation, typography, and sentence structure while preserving mathematical content. This supplied editorial label does not resolve the formulation gaps or constitute mathematical verification of this draft.
\end{itemize}
\end{document}
